\section{Background and related work}
\label{sec:background}

Every structural idea in \openjev{} exists in the literature. We state the
provenance plainly, because the paper's contribution is empirical rather
than architectural.

\paragraph{Scoring options instead of generating them.}
UniMC~\cite{unimc} recasts multiple-choice NLU by inserting an
option-marker token before each candidate and reading a decision from the
encoder's state at that position, avoiding generation entirely. Our readout
is the same idea. The consequence people usually reach for, that such a
model ``cannot hallucinate'', is narrower than it sounds: the model
cannot emit a string outside the menu, which says nothing about whether the
option it picks is right.

\paragraph{Encoding a shared prefix once.}
Parallel Context Windows~\cite{pcw} and related work formalise attention
masks in which several segments each attend to a common prefix and to
themselves, but not to one another. We use exactly this to isolate
questions, so that an answer does not depend on which other questions were
asked in the same call. On a causal backbone the prefix is ordinary
left-context, which is what a KV cache is for.

\paragraph{Multi-task training for zero-shot transfer.}
T0~\cite{t0} and the Flan collection~\cite{flan} establish that training
across many tasks produces zero-shot ability on unseen ones, and that the
gain continues accruing with task count into the high hundreds. Our mixture
is assembled from \texttt{tasksource}~\cite{tasksource}, which normalises a
large number of classification and multiple-choice datasets into a common
interface. Flan's task-count finding is the reason we treat 279 tasks as a
floor rather than a target.

\paragraph{Intermediate-task transfer.}
Training on an intermediate task before the target is
STILTs~\cite{stilts}; the domain-adaptive variant is
Gururangan et al.~\cite{dontstoppretraining}. Both report that the gain is largest when the
target dataset is small, and that it is inconsistent across task pairs and
can be negative. We observe both effects: a large gain at 395 target
examples (\S\ref{sec:res-router}), and two mixture changes that made
held-out transfer \emph{worse} (\S\ref{sec:ablations}).

\paragraph{Parameter-efficient adaptation.}
LoRA~\cite{lora} trains low-rank updates beside frozen weights. Our use is
conventional; the finding that it \emph{outperforms} full fine-tuning here,
rather than merely approximating it cheaply, is reported in
\S\ref{sec:res-router} together with the most likely explanation.

\paragraph{Ordinal and rejection heads.}
CORN/CORAL~\cite{corn} give rank-consistent ordinal objectives, and
adaptive decision boundaries~\cite{adb} detect out-of-scope inputs without
negative samples. We implement neither; \score{} is currently a softmax
over levels whose order is used only in rendering and in the
expectation-based readout. This is a limitation and we mark it as such
(\S\ref{sec:discussion}).

\paragraph{The system we compare against.}
Our reference is a commercial typed-decision API, \texttt{jev-1.13.0} from
TypeSafe, measured by black-box probing on 20 September 2026 at the sample
sizes stated. We inferred its design from behaviour: latency is flat from 2
to 255 options and an eight-question call costs the same wall time as a
one-question call, which is the signature of a shared prefix with parallel
readout. All figures attributed to it are one-day measurements of a hosted
endpoint from one network location, and should be reproduced before being
relied on.

Laya, from Convai Innovations, is the closest open-source system: the
same three primitives behind a near-identical
interface.\footnote{\url{https://huggingface.co/convaiinnovations/laya}}
We did not read its modelling code while designing ours, so that
converging on a similar design would be evidence rather than imitation,
and we score it in \S\ref{sec:res-laya} rather than only citing it.
